% Nguồn: Tuyển tập đề tham khảo môn Toán TS10 2025-2026 – TP. Hồ Chí Minh
% Đợt 9: Quận 8 (Đề 2) và Quận 10 (Đề 1) (14 câu tự luận)

% ==============================================================================
% QUẬN 8 – ĐỀ 2
% ==============================================================================

\begin{bt}%[9D6H18-2]
% nguon: Đề tham khảo Tuyển sinh 10 TP.HCM 2025-2026 – Quận 8 – Đề 2 – Bài 1
Cho hàm số $y = -\dfrac{1}{4}x^2$ có đồ thị là $(P)$.
\begin{enumerate}[a)]
    \item Vẽ $(P)$ trên hệ trục tọa độ $Oxy$.
    \item Tìm các điểm trên $(P)$ có tung độ bằng $-36$.
\end{enumerate}
\loigiai{
\begin{enumerate}[a)]
    \item Bảng giá trị tương ứng của $x$ và $y$:
    \begin{center}
    \begin{tabular}{|c|c|c|c|c|c|}
        \hline
        $x$ & $-4$ & $-2$ & $0$ & $2$ & $4$ \\ \hline
        $y = -\dfrac{1}{4}x^2$ & $-4$ & $-1$ & $0$ & $-1$ & $-4$ \\ \hline
    \end{tabular}
    \end{center}
    Đồ thị $(P)$ là một parabol đi qua gốc tọa độ $O(0;0)$, nhận trục tung $Oy$ làm trục đối xứng và bề lõm quay xuống dưới.
    \item Điểm $M(x; y) \in (P)$ có tung độ $y = -36$.\\
    Thay $y = -36$ vào phương trình parabol $(P)$, ta được:
    $-36 = -\dfrac{1}{4}x^2 \Leftrightarrow x^2 = 144 \Leftrightarrow x = \pm 12$.\\
    Vậy có hai điểm thuộc $(P)$ có tung độ bằng $-36$ là $(12; -36)$ và $(-12; -36)$.
\end{enumerate}
}
\end{bt}

\begin{bt}%[9D6H20-2]
% nguon: Đề tham khảo Tuyển sinh 10 TP.HCM 2025-2026 – Quận 8 – Đề 2 – Bài 2
Cho phương trình $2x^2 + x - 5 = 0$ \quad (1).
\begin{enumerate}[a)]
    \item Chứng tỏ phương trình (1) có hai nghiệm phân biệt.
    \item Giả sử $x_1, x_2$ là nghiệm của phương trình (1). Không giải phương trình, hãy tính giá trị của biểu thức: $A = x_1^2 + x_2^2 - x_1 - x_2$.
\end{enumerate}
\loigiai{
\begin{enumerate}[a)]
    \item Phương trình (1) có các hệ số $a = 2$ và $c = -5$.\\
    Vì $a \cdot c = 2 \cdot (-5) = -10 < 0$ nên phương trình (1) luôn có hai nghiệm phân biệt $x_1, x_2$.
    \item Theo định lí Viète, ta có:
    $\left\{\begin{aligned} &S = x_1 + x_2 = -\dfrac{1}{2} \\ &P = x_1 x_2 = -\dfrac{5}{2}. \end{aligned}\right.$\\
    Biến đổi biểu thức $A$:
    $\begin{aligned}
    A &= (x_1^2 + x_2^2) - (x_1 + x_2) = \left[(x_1 + x_2)^2 - 2x_1 x_2\right] - (x_1 + x_2) \\
    &= S^2 - 2P - S.
    \end{aligned}$\\
    Thay $S = -\dfrac{1}{2}$ và $P = -\dfrac{5}{2}$ vào biểu thức, ta được:
    $A = \left(-\dfrac{1}{2}\right)^2 - 2 \cdot \left(-\dfrac{5}{2}\right) - \left(-\dfrac{1}{2}\right) = \dfrac{1}{4} + 5 + \dfrac{1}{2} = \dfrac{23}{4} = 5{,}75$.\\
    Vậy $A = \dfrac{23}{4} = 5{,}75$.
\end{enumerate}
}
\end{bt}

\begin{bt}%[9D8H26-3]
% nguon: Đề tham khảo Tuyển sinh 10 TP.HCM 2025-2026 – Quận 8 – Đề 2 – Bài 3
Biểu đồ cột kép thống kê số lượng học sinh giỏi và học sinh khá của lớp 9A (lớp chỉ gồm học sinh giỏi và khá):
\begin{itemize}
    \item Tổ 1: 5 học sinh giỏi, 6 học sinh khá.
    \item Tổ 2: 6 học sinh giỏi, 5 học sinh khá.
    \item Tổ 3: 4 học sinh giỏi, 7 học sinh khá.
    \item Tổ 4: 5 học sinh giỏi, 6 học sinh khá.
\end{itemize}
\begin{enumerate}[a)]
    \item Tính số học sinh của lớp 9A.
    \item Giáo viên chủ nhiệm chọn ra 1 học sinh làm lớp trưởng, sau đó chọn tiếp 1 học sinh làm lớp phó. Biết cả lớp trưởng và lớp phó được chọn đều phải là học sinh giỏi. Tính xác suất để chọn ra được 1 bạn lớp phó (sau khi đã chọn xong bạn lớp trưởng).
\end{enumerate}
\loigiai{
\begin{enumerate}[a)]
    \item Tổng số học sinh của lớp 9A là:
    $(5 + 6) + (6 + 5) + (4 + 7) + (5 + 6) = 11 + 11 + 11 + 11 = 44$ (học sinh).\\
    Trong đó:
    \begin{itemize}
        \item Tổng số học sinh giỏi là: $5 + 6 + 4 + 5 = 20$ (học sinh).
        \item Tổng số học sinh khá là: $6 + 5 + 7 + 6 = 24$ (học sinh).
    \end{itemize}
    \item Ban đầu có 20 học sinh giỏi trong tổng số 44 học sinh của lớp.\\
    Sau khi đã chọn 1 bạn học sinh giỏi làm lớp trưởng thì trong lớp còn lại:
    $44 - 1 = 43$ (học sinh),\\
    và số học sinh giỏi còn lại là:
    $20 - 1 = 19$ (học sinh giỏi).\\
    Xác suất để chọn được 1 bạn làm lớp phó là học sinh giỏi từ các bạn còn lại là:
    $P = \dfrac{19}{43}$.
\end{enumerate}
}
\end{bt}

\begin{bt}%[9D6V21-3]
% nguon: Đề tham khảo Tuyển sinh 10 TP.HCM 2025-2026 – Quận 8 – Đề 2 – Bài 4
Trên một miếng đất hình chữ nhật có chiều dài là $(2x+10)\text{ m}$ và chiều rộng là $x\text{ m}$ ($x > 0$), người ta làm một lối đi hình bình hành có đáy nằm trên chiều dài với bề rộng là $2\text{ m}$. Phần còn lại là hai miếng đất hình thang vuông có diện tích bằng nhau dùng để trồng hoa hồng và hoa cúc.
\begin{enumerate}[a)]
    \item Viết biểu thức $A$ biểu diễn theo $x$ diện tích phần đất trồng hoa hồng.
    \item Người ta dự kiến lát sỏi lối đi với chi phí mỗi mét vuông là $120$ nghìn đồng. Hỏi tổng chi phí làm lối đi là bao nhiêu? Biết diện tích miếng đất trồng hoa hồng là $45\text{ m}^2$.
\end{enumerate}
\loigiai{
\begin{enumerate}[a)]
    \item Tổng diện tích miếng đất hình chữ nhật ban đầu là:
    $S_{\text{đất}} = x(2x + 10) = 2x^2 + 10x\text{ (m}^2\text{)}$.\\
    Lối đi hình bình hành có đáy bằng $2\text{ m}$ và chiều cao tương ứng chính là chiều rộng $x\text{ (m)}$ của mảnh đất hình chữ nhật.\\
    Diện tích lối đi là:
    $S_{\text{lối đi}} = 2 \cdot x = 2x\text{ (m}^2\text{)}$.\\
    Tổng diện tích đất còn lại để trồng hoa hồng và hoa cúc là:
    $S_{\text{còn}} = S_{\text{đất}} - S_{\text{lối đi}} = (2x^2 + 10x) - 2x = 2x^2 + 8x\text{ (m}^2\text{)}$.\\
    Vì hai phần đất trồng hoa hồng và hoa cúc có diện tích bằng nhau nên diện tích trồng hoa hồng là:
    $A = \dfrac{2x^2 + 8x}{2} = x^2 + 4x\text{ (m}^2\text{)}$.
    \item Theo đề bài, diện tích trồng hoa hồng là $45\text{ m}^2$, ta có phương trình:
    $x^2 + 4x = 45 \Leftrightarrow x^2 + 4x - 45 = 0 \Leftrightarrow (x - 5)(x + 9) = 0 \Leftrightarrow \left[\begin{aligned} &x = 5\text{ (thỏa mãn } x > 0) \\ &x = -9\text{ (loại)}. \end{aligned}\right.$\\
    Với $x = 5\text{ m}$, diện tích lối đi là:
    $S_{\text{lối đi}} = 2 \cdot 5 = 10\text{ (m}^2\text{)}$.\\
    Chi phí để lát sỏi làm lối đi là:
    $10 \cdot 120\,000 = 1\,200\,000\text{ (đồng)}$.\\
    Vậy chi phí để làm lối đi là $1\,200\,000$ đồng ($1{,}2$ triệu đồng).
\end{enumerate}
}
\end{bt}

\begin{bt}%[9H0V31-3]
% nguon: Đề tham khảo Tuyển sinh 10 TP.HCM 2025-2026 – Quận 8 – Đề 2 – Bài 5
Một xe chở hàng có thùng xe phía sau là một hình hộp chữ nhật có kích thước bên trong thùng xe: chiều dài $2{,}5\text{ m}$, chiều rộng $2\text{ m}$ và chiều cao $1{,}5\text{ m}$ (chiều dày thành xe không đáng kể).
\begin{enumerate}[a)]
    \item Tính thể tích của thùng chứa hàng.
    \item Một kho hàng có 220 thùng hàng dạng hình lập phương cạnh $50\text{ cm}$. Xe cần chuyển hết các thùng hàng đi nơi khác. Hỏi xe cần chở ít nhất bao nhiêu chuyến để chở hết 220 thùng hàng đó?
\end{enumerate}
\loigiai{
\begin{enumerate}[a)]
    \item Thể tích của thùng chứa hàng là:
    $V = 2{,}5 \cdot 2 \cdot 1{,}5 = 7{,}5\text{ (m}^3\text{)}$.
    \item Đổi cạnh mỗi thùng hàng dạng hình lập phương: $50\text{ cm} = 0{,}5\text{ m}$.\\
    Khi xếp các thùng hàng vào thùng xe tải:
    \begin{itemize}
        \item Số thùng hàng xếp được theo chiều dài là: $\dfrac{2{,}5}{0{,}5} = 5$ (thùng).
        \item Số thùng hàng xếp được theo chiều rộng là: $\dfrac{2}{0{,}5} = 4$ (thùng).
        \item Số thùng hàng xếp được theo chiều cao là: $\dfrac{1{,}5}{0{,}5} = 3$ (thùng).
    \end{itemize}
    Số thùng hàng tối đa xe có thể chở được trong một chuyến là:
    $5 \cdot 4 \cdot 3 = 60\text{ (thùng)}$.\\
    Số chuyến xe ít nhất cần chở để vận chuyển hết 220 thùng hàng là:
    $N = \left\lceil \dfrac{220}{60} \right\rceil = \lceil 3{,}67 \rceil = 4\text{ (chuyến)}$.\\
    Vậy xe cần chở ít nhất 4 chuyến.
\end{enumerate}
}
\end{bt}

\begin{bt}%[9D1H3-4]
% nguon: Đề tham khảo Tuyển sinh 10 TP.HCM 2025-2026 – Quận 8 – Đề 2 – Bài 6
Đồng thau là hợp kim có chứa đồng và kẽm. Trong đồng thau màu đỏ có khối lượng kẽm chiếm $20\%$, trong đồng thau màu vàng nhạt có khối lượng kẽm chiếm $42\%$. Từ một ống đồng thau màu đỏ có khối lượng $500\text{ g}$, người ta cần cho thêm bao nhiêu gam kẽm để tạo ra một ống đồng thau màu vàng nhạt? (Làm tròn kết quả đến hàng đơn vị).
\loigiai{
Khối lượng kẽm có sẵn trong $500\text{ g}$ đồng thau màu đỏ là:
$500 \cdot 20\% = 100\text{ (g)}$.\\
Gọi $x$ (g) là khối lượng kẽm cần cho thêm vào ($x > 0$).\\
Khối lượng của hợp kim đồng thau mới sau khi thêm kẽm là:
$500 + x\text{ (g)}$.\\
Khối lượng kẽm trong hợp kim mới là:
$100 + x\text{ (g)}$.\\
Vì hợp kim mới là đồng thau màu vàng nhạt chứa $42\%$ kẽm nên ta có phương trình:
$\dfrac{100 + x}{500 + x} = 42\% \Leftrightarrow \dfrac{100 + x}{500 + x} = 0{,}42 \Leftrightarrow 100 + x = 0{,}42(500 + x) \Leftrightarrow 100 + x = 210 + 0{,}42x \Leftrightarrow 0{,}58x = 110 \Leftrightarrow x = \dfrac{110}{0{,}58} \approx 189{,}66\text{ (g)}$.\\
Làm tròn kết quả đến hàng đơn vị, ta được $x \approx 190\text{ g}$.\\
Vậy cần cho thêm khoảng $190\text{ g}$ kẽm.
}
\end{bt}

\begin{bt}%[9H9C29-3]
% nguon: Đề tham khảo Tuyển sinh 10 TP.HCM 2025-2026 – Quận 8 – Đề 2 – Bài 7
Cho đường tròn $(O, R)$ đường kính $AB$. Vẽ hai tiếp tuyến $Ax$ và $By$. Điểm $E$ thuộc tia $BA$ sao cho $A$ là trung điểm của $EO$. Từ $E$ vẽ tiếp tuyến $EM$ ($M$ là tiếp điểm) cắt $Ax$ tại $C$ và $By$ tại $D$.
\begin{enumerate}[a)]
    \item Chứng minh tứ giác $OMDB$ là tứ giác nội tiếp.
    \item Đoạn thẳng $AM$ cắt $OC$ tại $I$, đoạn thẳng $MB$ cắt $OD$ tại $K$. Chứng minh tứ giác $OKMI$ là hình chữ nhật và tam giác $OKI$ đồng dạng với tam giác $OCD$.
    \item Tính diện tích tứ giác $KICD$ theo $R$.
\end{enumerate}
\loigiai{
\begin{enumerate}[a)]
    \item Vì $By$ là tiếp tuyến tại $B$ và $EM$ là tiếp tuyến tại $M$ của đường tròn $(O)$ nên:
    $\widehat{OBD} = 90^\circ$ và $\widehat{OMD} = 90^\circ$.\\
    Tứ giác $OMDB$ có hai đỉnh kề $B$ và $M$ cùng nhìn cạnh $OD$ dưới một góc vuông, hay $\widehat{OBD} = \widehat{OMD} = 90^\circ$.\\
    Vậy tứ giác $OMDB$ nội tiếp đường tròn đường kính $OD$.
    \item Theo tính chất hai tiếp tuyến cắt nhau từ $C$ đến $(O)$ (tiếp tuyến $CA$ và $CM$), ta có $CA = CM$ và $OA = OM = R$, suy ra $OC$ là đường trung trực của đoạn $AM$, do đó $OC \perp AM$ tại $I$.\\
    Tương tự, $OD$ là đường trung trực của đoạn thẳng $BM$, do đó $OD \perp BM$ tại $K$.\\
    Vì điểm $M$ thuộc đường tròn đường kính $AB$ nên $\widehat{AMB} = 90^\circ$ (góc nội tiếp chắn nửa đường tròn).\\
    Tứ giác $OKMI$ có 3 góc vuông: $\widehat{OIM} = 90^\circ$, $\widehat{OKM} = 90^\circ$ và $\widehat{IMK} = 90^\circ$.\\
    Vậy tứ giác $OKMI$ là hình chữ nhật.\\
    Vì $OKMI$ là hình chữ nhật nên $\widehat{IOK} = 90^\circ$ và hai đường chéo bằng nhau: $IK = OM = R$.\\
    Xét tam giác $OCD$ vuông tại $O$ (vì $OC$ và $OD$ là hai tia phân giác của hai góc kề bù $\widehat{AOM}$ và $\widehat{BOM}$):\\
    Trong tam giác vuông $OCD$ có $OM \perp CD$ tại $M$, áp dụng hệ thức lượng ta có:
    $OI \cdot OC = OM^2 = OK \cdot OD \Rightarrow \dfrac{OI}{OD} = \dfrac{OK}{OC}$.\\
    Xét $\triangle OKI$ và $\triangle OCD$ có:
    $\left\{\begin{aligned} &\widehat{COD}\text{ là góc chung} \\ &\dfrac{OI}{OD} = \dfrac{OK}{OC}\text{ (chứng minh trên)}. \end{aligned}\right.$\\
    Do đó $\triangle OKI \sim \triangle OCD$ (c-g-c).
    \item Vì $\triangle OKI \sim \triangle OCD$ với tỉ số đồng dạng $k = \dfrac{IK}{CD}$:\\
    Ta có $A$ là trung điểm của $EO$ nên $OA = R \Rightarrow EO = 2R$.\\
    Xét tam giác $EAC$ vuông tại $A$ và tam giác $EMO$ vuông tại $M$ có góc $\widehat{OEM}$ chung, do đó $\triangle EAC \sim \triangle EMO$ (g-g).\\
    Suy ra $\dfrac{AC}{MO} = \dfrac{EA}{EM}$.\\
    Tam giác $EMO$ vuông tại $M$ có $EO = 2R, OM = R \Rightarrow EM = \sqrt{EO^2 - OM^2} = \sqrt{4R^2 - R^2} = R\sqrt{3}$.\\
    Do đó $AC = MO \cdot \dfrac{EA}{EM} = R \cdot \dfrac{R}{R\sqrt{3}} = \dfrac{R}{\sqrt{3}}$.\\
    Trong tam giác vuông $OCD$ có đường cao $OM$:
    $OM^2 = AC \cdot BD \Rightarrow BD = \dfrac{OM^2}{AC} = \dfrac{R^2}{\dfrac{R}{\sqrt{3}}} = R\sqrt{3}$.\\
    Độ dài đoạn thẳng $CD = CM + MD = AC + BD = \dfrac{R}{\sqrt{3}} + R\sqrt{3} = \dfrac{4R}{\sqrt{3}}$.\\
    Diện tích tam giác vuông $OCD$ là:
    $S_{\triangle OCD} = \dfrac{1}{2} OM \cdot CD = \dfrac{1}{2} \cdot R \cdot \dfrac{4R}{\sqrt{3}} = \dfrac{2R^2}{\sqrt{3}} = \dfrac{2\sqrt{3}R^2}{3}$.\\
    Tỉ số đồng dạng giữa $\triangle OKI$ và $\triangle OCD$ là:
    $k = \dfrac{IK}{CD} = \dfrac{R}{\dfrac{4R}{\sqrt{3}}} = \dfrac{\sqrt{3}}{4}$.\\
    Do đó tỉ số diện tích là:
    $\dfrac{S_{\triangle OKI}}{S_{\triangle OCD}} = k^2 = \left(\dfrac{\sqrt{3}}{4}\right)^2 = \dfrac{3}{16}$.\\
    Diện tích tứ giác $KICD$ bằng hiệu diện tích hai tam giác:
    $S_{KICD} = S_{\triangle OCD} - S_{\triangle OKI} = S_{\triangle OCD} \cdot \left(1 - \dfrac{3}{16}\right) = \dfrac{13}{16} S_{\triangle OCD} = \dfrac{13}{16} \cdot \dfrac{2\sqrt{3}R^2}{3} = \dfrac{13\sqrt{3}R^2}{24}$.\\
    Vậy diện tích tứ giác $KICD$ là $\dfrac{13\sqrt{3}R^2}{24}$.
\end{enumerate}
}
\end{bt}

% ==============================================================================
% QUẬN 10 – ĐỀ 1
% ==============================================================================

\begin{bt}%[9D6H18-2]
% nguon: Đề tham khảo Tuyển sinh 10 TP.HCM 2025-2026 – Quận 10 – Đề 1 – Bài 1
Cho parabol $(P)\colon y = \dfrac{1}{2}x^2$.
\begin{enumerate}[a)]
    \item Vẽ đồ thị $(P)$ trên hệ trục tọa độ $Oxy$.
    \item Tìm điểm $M$ thuộc $(P)$ có tung độ gấp đôi hoành độ và khác gốc tọa độ $O$.
\end{enumerate}
\loigiai{
\begin{enumerate}[a)]
    \item Bảng giá trị tương ứng của $x$ và $y$:
    \begin{center}
    \begin{tabular}{|c|c|c|c|c|c|}
        \hline
        $x$ & $-4$ & $-2$ & $0$ & $2$ & $4$ \\ \hline
        $y = \dfrac{1}{2}x^2$ & $8$ & $2$ & $0$ & $2$ & $8$ \\ \hline
    \end{tabular}
    \end{center}
    Đồ thị $(P)$ là một parabol đi qua gốc tọa độ $O(0;0)$, nhận trục tung $Oy$ làm trục đối xứng và nằm phía trên trục hoành.
    \item Điểm $M(x; y) \in (P)$ có tung độ gấp đôi hoành độ, tức là $y = 2x$.\\
    Vì $M \in (P)$ nên ta có phương trình:
    $\dfrac{1}{2}x^2 = 2x \Leftrightarrow x^2 - 4x = 0 \Leftrightarrow x(x - 4) = 0 \Leftrightarrow \left[\begin{aligned} &x = 0 \\ &x = 4. \end{aligned}\right.$\\
    Vì điểm $M$ khác gốc tọa độ $O(0;0)$ nên $x = 4 \Rightarrow y = 2 \cdot 4 = 8$.\\
    Vậy tọa độ điểm cần tìm là $M(4; 8)$.
\end{enumerate}
}
\end{bt}

\begin{bt}%[9D6H20-2]
% nguon: Đề tham khảo Tuyển sinh 10 TP.HCM 2025-2026 – Quận 10 – Đề 1 – Bài 2
Cho phương trình $x^2 - 4x - 1 = 0$.
\begin{enumerate}[a)]
    \item Chứng minh phương trình trên có hai nghiệm phân biệt.
    \item Không giải phương trình, hãy tính giá trị của biểu thức: $A = (x_1 - 2x_2)(x_2 - x_1) - x_1(x_1 + 2025x_2)$.
\end{enumerate}
\loigiai{
\begin{enumerate}[a)]
    \item Phương trình $x^2 - 4x - 1 = 0$ có các hệ số $a = 1$ và $c = -1$.\\
    Vì $a \cdot c = 1 \cdot (-1) = -1 < 0$ nên phương trình luôn có hai nghiệm phân biệt $x_1, x_2$.
    \item Theo định lí Viète, ta có:
    $\left\{\begin{aligned} &S = x_1 + x_2 = 4 \\ &P = x_1 x_2 = -1. \end{aligned}\right.$\\
    Khai triển biểu thức $A$:
    $\begin{aligned}
    A &= (x_1 - 2x_2)(x_2 - x_1) - x_1(x_1 + 2025x_2) \\
    &= (x_1 x_2 - x_1^2 - 2x_2^2 + 2x_1 x_2) - (x_1^2 + 2025x_1 x_2) \\
    &= -2x_1^2 - 2x_2^2 - 2022x_1 x_2 \\
    &= -2(x_1^2 + x_2^2) - 2022x_1 x_2 \\
    &= -2\left[(x_1 + x_2)^2 - 2x_1 x_2\right] - 2022x_1 x_2 \\
    &= -2(S^2 - 2P) - 2022P.
    \end{aligned}$\\
    Thay $S = 4$ và $P = -1$ vào biểu thức:
    $A = -2\left[4^2 - 2(-1)\right] - 2022(-1) = -2(16 + 2) + 2022 = -36 + 2022 = 1986$.\\
    Vậy $A = 1986$.
\end{enumerate}
}
\end{bt}

\begin{bt}%[9D1V3-4]
% nguon: Đề tham khảo Tuyển sinh 10 TP.HCM 2025-2026 – Quận 10 – Đề 1 – Bài 3
Một cửa hàng bán sỉ nước giải khát Pepsi với giá niêm yết $182\,000$ đồng/thùng. Mỗi tháng cửa hàng xuất bán trung bình được $550$ thùng. Để mở rộng kinh doanh, cửa hàng đã giảm giá $x\%$ cho mỗi thùng Pepsi và đã bán thêm được $y$ thùng trong một tháng.
\begin{enumerate}[a)]
    \item Viết biểu thức số tiền thu được trong một tháng sau khi mở rộng kinh doanh theo $x$ và $y$.
    \item Biết trong tháng đầu tiên sau khi mở rộng, cửa hàng giảm giá $5\%$ và thu được số tiền nhiều hơn khi chưa mở rộng là $38{,}22$ triệu đồng. Hãy tính xem cửa hàng đã bán thêm được bao nhiêu thùng Pepsi trong tháng đó.
\end{enumerate}
\loigiai{
\begin{enumerate}[a)]
    \item Giá bán của mỗi thùng Pepsi sau khi giảm giá $x\%$ là:
    $182\,000 \cdot (1 - x\%)\text{ (đồng)}$.\\
    Số lượng thùng Pepsi bán ra trong một tháng sau khi mở rộng là:
    $550 + y\text{ (thùng)}$.\\
    Biểu thức số tiền thu được trong một tháng theo $x$ và $y$ là:
    $T(x; y) = 182\,000(1 - x\%)(550 + y)\text{ (đồng)}$.
    \item Số tiền cửa hàng thu được mỗi tháng khi chưa mở rộng kinh doanh là:
    $182\,000 \cdot 550 = 100\,100\,000\text{ (đồng)} = 100{,}1\text{ (triệu đồng)}$.\\
    Trong tháng đầu tiên sau khi giảm giá $5\%$, số tiền thu được là:
    $100{,}1 + 38{,}22 = 138{,}32\text{ (triệu đồng)} = 138\,320\,000\text{ (đồng)}$.\\
    Giá bán của mỗi thùng Pepsi sau khi giảm giá $5\%$ là:
    $182\,000 \cdot (1 - 5\%) = 182\,000 \cdot 0{,}95 = 172\,900\text{ (đồng)}$.\\
    Tổng số thùng Pepsi bán ra trong tháng này là:
    $\dfrac{138\,320\,000}{172\,900} = 800\text{ (thùng)}$.\\
    Số thùng Pepsi cửa hàng đã bán thêm được là:
    $y = 800 - 550 = 250\text{ (thùng)}$.\\
    Vậy cửa hàng đã bán thêm được 250 thùng Pepsi.
\end{enumerate}
}
\end{bt}

\begin{bt}%[9H0V31-3]
% nguon: Đề tham khảo Tuyển sinh 10 TP.HCM 2025-2026 – Quận 10 – Đề 1 – Bài 4
Một cây bút chì hình trụ có chiều dài $180\text{ mm}$ và đường kính $7{,}2\text{ mm}$. Phần ruột bút làm bằng chì hình trụ có cùng chiều dài với bút và đường kính ngòi bằng $3{,}4\text{ mm}$.
\begin{enumerate}[a)]
    \item Tính thể tích chì cần dùng để làm lõi một cây bút chì khi chưa gọt (làm tròn kết quả đến hàng phần trăm).
    \item Để có phần vỏ gỗ của bút chì, người ta dùng các thanh gỗ hình hộp chữ nhật có đáy là hình vuông cạnh $8\text{ mm}$ và chiều dài $185\text{ mm}$. Hỏi với $10\text{ m}^3$ gỗ chuyên dụng thì có thể tạo ra được bao nhiêu cây bút chì, biết phần hao hụt khi xẻ nhỏ gỗ chiếm $12\%$ tổng lượng gỗ đem làm?
\end{enumerate}
\loigiai{
\begin{enumerate}[a)]
    \item Bán kính của phần ruột chì hình trụ là:
    $r = \dfrac{3{,}4}{2} = 1{,}7\text{ (mm)}$.\\
    Chiều dài của lõi chì là $h = 180\text{ mm}$.\\
    Thể tích chì cần dùng để làm lõi một cây bút chì là:
    $V = \pi r^2 h = \pi \cdot (1{,}7)^2 \cdot 180 = 520{,}2\pi \approx 1634{,}26\text{ (mm}^3\text{)}$.
    \item Thể tích của một thanh gỗ hình hộp chữ nhật làm vỏ bút chì là:
    $V_1 = 8 \cdot 8 \cdot 185 = 11\,840\text{ (mm}^3\text{)}$.\\
    Đổi thể tích khối gỗ lớn ra $\text{mm}^3$:
    $10\text{ m}^3 = 10 \cdot 10^9 = 10^{10}\text{ (mm}^3\text{)}$.\\
    Vì hao hụt chiếm $12\%$ nên thể tích gỗ thực tế được đưa vào sản xuất là:
    $V_{\text{thực}} = 10^{10} \cdot (1 - 12\%) = 8{,}8 \cdot 10^9\text{ (mm}^3\text{)}$.\\
    Số lượng vỏ bút chì (tương ứng với số cây bút chì) có thể sản xuất được là:
    $N = \left\lfloor \dfrac{8{,}8 \cdot 10^9}{11\,840} \right\rfloor \approx \lfloor 743\,243{,}24 \rfloor = 743\,243\text{ (cây bút)}$.\\
    Vậy có thể tạo ra được $743\,243$ cây bút chì.
\end{enumerate}
}
\end{bt}

\begin{bt}%[9D1H3-4]
% nguon: Đề tham khảo Tuyển sinh 10 TP.HCM 2025-2026 – Quận 10 – Đề 1 – Bài 5
Một công ty phân bón cần sản xuất một loại phân bón chứa $30\%$ potassium từ hai loại nguyên liệu: loại A chứa $24\%$ potassium và loại B chứa $40\%$ potassium.
\begin{enumerate}[a)]
    \item Tính khối lượng mỗi loại nguyên liệu cần sử dụng để thu được hỗn hợp $500\text{ kg}$ chứa $30\%$ potassium.
    \item Nếu công ty thêm $50\text{ kg}$ loại A vào hỗn hợp $500\text{ kg}$ ban đầu, tỉ lệ phần trăm potassium trong hỗn hợp mới là bao nhiêu (làm tròn số đến hàng phần mười)?
\end{enumerate}
\loigiai{
\begin{enumerate}[a)]
    \item Gọi $x, y$ lần lượt là khối lượng nguyên liệu loại A và loại B cần dùng ($x, y > 0$, đơn vị kg).\\
    Tổng khối lượng hỗn hợp cần sản xuất là $500\text{ kg}$, ta có:
    $x + y = 500$.\\
    Khối lượng potassium nguyên chất trong $500\text{ kg}$ hỗn hợp chứa $30\%$ là:
    $500 \cdot 30\% = 150\text{ (kg)}$.\\
    Theo bài ra ta có phương trình:
    $24\% x + 40\% y = 150 \Leftrightarrow 0{,}24x + 0{,}4y = 150$.\\
    Ta có hệ phương trình:
    $\left\{\begin{aligned} &x + y = 500 \\ &0{,}24x + 0{,}4y = 150 \end{aligned}\right. \Leftrightarrow \left\{\begin{aligned} &0{,}4x + 0{,}4y = 200 \\ &0{,}24x + 0{,}4y = 150 \end{aligned}\right. \Leftrightarrow \left\{\begin{aligned} &0{,}16x = 50 \\ &y = 500 - x \end{aligned}\right. \Leftrightarrow \left\{\begin{aligned} &x = 312{,}5 \\ &y = 187{,}5. \end{aligned}\right.$\\
    Vậy cần dùng $312{,}5\text{ kg}$ loại A và $187{,}5\text{ kg}$ loại B.
    \item Khi thêm $50\text{ kg}$ loại A vào hỗn hợp ban đầu:\\
    Tổng khối lượng của hỗn hợp mới là:
    $500 + 50 = 550\text{ (kg)}$.\\
    Khối lượng potassium trong $50\text{ kg}$ loại A thêm vào là:
    $50 \cdot 24\% = 12\text{ (kg)}$.\\
    Tổng khối lượng potassium trong hỗn hợp mới là:
    $150 + 12 = 162\text{ (kg)}$.\\
    Tỉ lệ phần trăm potassium trong hỗn hợp mới là:
    $\dfrac{162}{550} \cdot 100\% \approx 29{,}45\% \approx 29{,}5\%$.\\
    Vậy hàm lượng potassium trong hỗn hợp mới là khoảng $29{,}5\%$.
\end{enumerate}
}
\end{bt}

\begin{bt}%[9D8H26-3]
% nguon: Đề tham khảo Tuyển sinh 10 TP.HCM 2025-2026 – Quận 10 – Đề 1 – Bài 6
Bạn An tham gia trò chơi ném bi vào một bảng gồm 25 ô có các mức điểm: 5 ô có điểm 5; 12 ô có điểm 3; 4 ô có điểm -1; 4 ô có điểm -2. Ném ra ngoài bảng bị trừ 5 điểm.
Phần thưởng sau 10 lần ném: Đạt từ 30 điểm đến 49 điểm nhận phần quà trị giá 300.000 đồng.
\begin{enumerate}[a)]
    \item Trong 9 lần ném đầu, bạn An ném được 5 lần vào ô điểm 5, một lần ra ngoài bảng, 2 lần vào ô điểm 3 và một lần vào ô điểm (-1). Tính số điểm bạn An nhận được sau 9 lần ném.
    \item Bạn An có cơ hội nhận phần quà 300.000 đồng sau 10 lần ném không? Nếu có thì lần tiếp theo phải ném vào ô nào? Tính xác suất để bạn An nhận được phần quà đó biết lần thứ 10 bạn chắc chắn ném trúng vào trong bảng.
\end{enumerate}
\loigiai{
\begin{enumerate}[a)]
    \item Tổng số điểm bạn An đạt được sau 9 lần ném là:
    $5 \cdot 5 + 1 \cdot (-5) + 2 \cdot 3 + 1 \cdot (-1) = 25 - 5 + 6 - 1 = 25\text{ (điểm)}$.
    \item Để nhận được phần quà trị giá 300.000 đồng, tổng số điểm sau 10 lần ném phải đạt từ 30 đến 49 điểm.\\
    Do đó, số điểm ở lần ném thứ 10 cần đạt tối thiểu là:
    $30 - 25 = 5\text{ (điểm)}$.\\
    Trong bảng 25 ô, mức điểm cao nhất của một ô là 5 điểm.\\
    Do đó, bạn An vẫn có cơ hội nhận được phần quà này nếu ở lần thứ 10 bạn ném trúng vào một ô có điểm 5.\\
    Vì lần thứ 10 An chắc chắn ném trúng vào trong bảng nên số kết quả có thể xảy ra là 25 ô.\\
    Trên bảng có đúng 5 ô có điểm 5, nên số kết quả thuận lợi là 5.\\
    Xác suất để bạn An nhận được phần quà trị giá 300.000 đồng là:
    $P = \dfrac{5}{25} = \dfrac{1}{5} = 0{,}2$.
\end{enumerate}
}
\end{bt}

\begin{bt}%[9H9C29-3]
% nguon: Đề tham khảo Tuyển sinh 10 TP.HCM 2025-2026 – Quận 10 – Đề 1 – Bài 7
Cho tam giác $ABC$ có ba góc nhọn ($AB < AC$) nội tiếp đường tròn $(O; R)$. Ba đường cao $AD, BE, CF$ cắt nhau tại $H$. Gọi $M, N$ lần lượt là giao điểm của $(O)$ với các tia $BE, AD$ ($M \ne B, N \ne A$).
\begin{enumerate}[a)]
    \item Chứng minh tứ giác $ABDE$ nội tiếp và xác định tâm $I$ của đường tròn này, từ đó suy ra $DE \parallel MN$.
    \item Kẻ đường kính $CK$ của $(O)$. Chứng minh tứ giác $AKBH$ là hình bình hành và suy ra 3 điểm $H, I, K$ thẳng hàng.
    \item Trong trường hợp $\widehat{BCA} = 60^\circ$. Chứng minh $DE = \dfrac{1}{2}AB$ và tính diện tích hình viên phân giới hạn bởi cung nhỏ $DE$ và dây cung $DE$ của đường tròn $(I)$ theo $R$.
\end{enumerate}
\loigiai{
\begin{enumerate}[a)]
    \item Vì $AD \perp BC$ và $BE \perp AC$ nên $\widehat{ADB} = \widehat{AEB} = 90^\circ$.\\
    Tứ giác $ABDE$ có hai đỉnh kề $D$ và $E$ cùng nhìn cạnh $AB$ dưới một góc vuông.\\
    Vậy tứ giác $ABDE$ nội tiếp đường tròn đường kính $AB$, tâm $I$ là trung điểm của cạnh $AB$.\\
    Trong đường tròn $(I)$, góc nội tiếp $\widehat{BDE} = \widehat{BAE}$ (cùng chắn cung $BE$).\\
    Trong đường tròn $(O)$, góc nội tiếp $\widehat{BMN} = \widehat{BAN}$ (cùng chắn cung $BN$).\\
    Vì $A, H, D, N$ thẳng hàng nên $\widehat{BAE} \equiv \widehat{BAN}$, suy ra $\widehat{BDE} = \widehat{BMN}$.\\
    Vì hai góc này ở vị trí đồng vị nên $DE \parallel MN$.
    \item Vì $CK$ là đường kính của $(O)$ nên:
    $\widehat{CAK} = 90^\circ \Rightarrow AK \perp AC$.\\
    Mà $BH \perp AC$ (do $BE$ là đường cao), suy ra $AK \parallel BH$.\\
    Tương tự: $\widehat{CBK} = 90^\circ \Rightarrow BK \perp BC$.\\
    Mà $AH \perp BC$ (do $AD$ là đường cao), suy ra $BK \parallel AH$.\\
    Tứ giác $AKBH$ có các cặp cạnh đối song song ($AK \parallel BH$ và $BK \parallel AH$) nên $AKBH$ là hình bình hành.\\
    Vì $I$ là trung điểm của đường chéo $AB$ trong hình bình hành $AKBH$ nên $I$ đồng thời là trung điểm của đường chéo $HK$.\\
    Vậy ba điểm $H, I, K$ thẳng hàng.
    \item Khi $\widehat{BCA} = 60^\circ$:\\
    Tam giác $CDE$ có $\widehat{CDE} = \widehat{BAC}$ (do tứ giác $ABDE$ nội tiếp), suy ra $\triangle CDE \sim \triangle CAB$ (g-g).\\
    Tỉ số đồng dạng là:
    $k = \dfrac{DE}{AB} = \dfrac{CD}{CA} = \cos \widehat{BCA} = \cos 60^\circ = \dfrac{1}{2} \Rightarrow DE = \dfrac{1}{2}AB$.\\
    Đường tròn $(I)$ ngoại tiếp tứ giác $ABDE$ có bán kính là:
    $r = \dfrac{AB}{2}$.\\
    Vì $DE = \dfrac{1}{2}AB = r$ nên tam giác $IDE$ là tam giác đều cạnh $r$, suy ra góc ở tâm $\widehat{DIE} = 60^\circ$.\\
    Theo định lí sin trong tam giác $ABC$ nội tiếp $(O; R)$:
    $AB = 2R \sin \widehat{BCA} = 2R \sin 60^\circ = 2R \cdot \dfrac{\sqrt{3}}{2} = R\sqrt{3} \Rightarrow r = \dfrac{R\sqrt{3}}{2}$.\\
    Diện tích hình quạt tròn $IDE$ ứng với góc ở tâm $60^\circ$ là:
    $S_{\text{quạt}} = \dfrac{\pi r^2 \cdot 60}{360} = \dfrac{\pi r^2}{6}$.\\
    Diện tích tam giác đều $IDE$ cạnh $r$ là:
    $S_{\triangle IDE} = \dfrac{r^2\sqrt{3}}{4}$.\\
    Diện tích hình viên phân giới hạn bởi cung nhỏ $DE$ và dây cung $DE$ là:
    $S = S_{\text{quạt}} - S_{\triangle IDE} = r^2\left(\dfrac{\pi}{6} - \dfrac{\sqrt{3}}{4}\right) = \dfrac{3R^2}{4}\left(\dfrac{\pi}{6} - \dfrac{\sqrt{3}}{4}\right) = R^2\left(\dfrac{\pi}{8} - \dfrac{3\sqrt{3}}{16}\right)$.\\
    Vậy diện tích hình viên phân là $R^2\left(\dfrac{\pi}{8} - \dfrac{3\sqrt{3}}{16}\right)$.
\end{enumerate}
}
\end{bt}
